% Lösungen Einheit 2 – Nr. 24–31
\einheitenkopf[e2]{Einheit 2 von 4 · Satz vom Nullprodukt}

\erg{24}{a) ja \quad b) nein \quad c) ja \quad d) nein \quad e) nein, links steht eine Differenz \quad f) ja}
\erg{25}{a) $x_1 = 2$, $x_2 = 5$ \quad b) $x_1 = 1$, $x_2 = 4$ \quad c) $x_1 = 6$, $x_2 = 3$ \quad d) $x_1 = 7$, $x_2 = 2$ \quad e) $x_1 = -4$, $x_2 = 1$ \quad f) $x_1 = 0$, $x_2 = -6$ \quad g) $x = 3$ (eine Lösung)}
\erg{26}{a) (wA) \quad b) $(-4)\cdot 3 = -12$, (fA) \quad c) $(-7)\cdot 2 = -14$, (wA) \quad d) $(-1)\cdot(-4) = 4$, (fA) \quad e) $(-4)\cdot(-3) = 12$, (wA) \quad f) $x = -2$, denn $(-2)\cdot 5 = -10$}
\erg{27}{a) $x\cdot(x + 4) = 0$; $x_1 = 0$, $x_2 = -4$ \quad b) $x_1 = 0$, $x_2 = -5$ \quad c) $x_1 = 0$, $x_2 = -2$ \quad d) $x\cdot(x - 6) = 0$; $x_1 = 0$, $x_2 = 6$ \quad e) $x\cdot(5x - 15) = 0$; $x_1 = 0$, $x_2 = 3$}
\erg{28}{a) $x^2 + 3x - 10 = 0$ \quad b) $x^2 - 4x + 3 = 0$ \quad c) $x^2 - 2x + x - 2 = 4$, also $x^2 - x - 6 = 0$}
\erg{29}{a) $x_1 = -2$, $x_2 = -3$ \quad b) $x_1 = 4$, $x_2 = -2$ \quad c) $4\cdot(x^2 + 2x - 3x - 6) = 4x^2 - 4x - 24$; $x - 3 = 0$ oder $x + 2 = 0$; $x_1 = 3$, $x_2 = -2$}
\erg{30}{a) Durch $x$ geteilt, dabei geht $x = 0$ verloren. Richtig: $x^2 - 6x = 0$, $x\cdot(x - 6) = 0$; $x_1 = 0$, $x_2 = 6$. \quad b) Vorzeichen beim Nullsetzen: $x + 5 = 0$ ergibt $x = -5$, $x - 2 = 0$ ergibt $x = 2$. \quad c) Der Satz gilt nur, wenn rechts null steht. Richtig: $x^2 + x - 12 = 0$ ordnen, dann mit der Formel aus Einheit 3: $x_1 = 3$, $x_2 = -4$.}
\erg{31}{a) Multipliziert man eine Zahl mit null, ist das Ergebnis immer null – egal, wie groß der andere Faktor ist. \quad b) $x$ kann null sein, und durch null darf man nicht teilen; die Lösung $x = 0$ ginge verloren ($x^2 - 3x = 0$ hat die Lösungen 0 und 3).}
